\documentclass[10pt,a4paper]{amsart}
\usepackage[margin=19mm]{geometry}
\usepackage{amsmath,amssymb,mathtools}
\usepackage[scr=rsfs,cal=euler]{mathalfa}
\usepackage{xfrac}
\usepackage[unicode,hidelinks,bookmarks=false]{hyperref}
\newcommand{\ie}{i.e.\ }
\newcommand{\cf}{cf.\ }
\newcommand{\resp}{respectively\ }
\newcommand{\Subsection}{\subsection}
\providecommand{\nobreakdash}{\nobreak\hbox{-}}
\setlength{\emergencystretch}{2em}
\multlinegap0pt
\usepackage{bm}
\usepackage{bbm}
\usepackage{dsfont}
\usepackage{xspace}
\usepackage{cancel}
\usepackage{mathtools}
\usepackage{amsthm}
\makeatletter
\@ifundefined{Def}{\newtheorem{Def}{Def}}{}
\@ifundefined{Lem}{\newtheorem{Lem}[Def]{Lemma}}{}
\@ifundefined{Prop}{\newtheorem{Prop}[Def]{Proposition}}{}
\makeatother
\title[Moduli stacks of Higgs bundles on stable curves]{Moduli stacks of Higgs bundles on stable curves}
\author{Oren Ben-Bassat, Sourav Das, Tony Pantev}
\date{Source: arXiv:2310.07778v3 (CC BY 4.0)}
\begin{document}
\maketitle

\noindent\emph{Source and licence.} Moduli stacks of Higgs bundles on stable curves. Version arXiv:2310.07778v3 (CC BY 4.0), distributed under the Creative Commons Attribution 4.0 International licence, as recorded in the arXiv licence metadata for this exact version and shown on the abstract page https://arxiv.org/abs/2310.07778v3. Full rights notice: licenses/arxiv-2310.07778-CC-BY-4.0.txt.\par\smallskip

\noindent\textit{Editorial scope.} Counted unit: the Prop whose source label is \texttt{rellog101} in the article (source0.tex lines 602--604, source label \texttt{rellog101}), delivered together with the article's own complete original proof of it (source0.tex lines 605--674, one \texttt{proof} environment of 70 lines). No mathematics is written, repaired or re-derived: statement and proof are the article's own text line for line. Required context retained uncounted: et1 (lines 510--516). Every definition, display and label of those lines is retained unchanged. Omitted from this package and not redistributed: 1--509, 517--601, 675--2156. Nothing inside a retained excerpt is omitted or altered: every retained body is delivered exactly as the article's lines are written, so no omission receipt of any kind accompanies this package. Citations: the retained excerpts carry no citation command at all, so no bibliography entry of the article is transcribed and no reference is redistributed. Reference adaptation: none was needed and none was made; every reference inside the delivered unit resolves inside it, so no unresolved reference or citation is produced. Printed slips kept literal: no printed word, symbol, hypothesis or number of the counted statement or of its proof was altered. Layout and class: the article's own class is \texttt{amsart} with packages including xy, amsmath, amssymb, amscd, amsthm, mathtools, amsfonts, mathrsfs, tikz-cd, bbm, enumitem, euscript, inputenc, babel; the wrapper is the lane's pinned \texttt{amsart} editorial wrapper at 10pt on A4, which restates the theorem-like environments the retained text needs and adds the article's own macro definitions that the retained text needs (none), with extra wrapper packages (bm, bbm, dsfont, xspace, cancel, mathtools). The delivered numbering is the unit's own. Assets: no retained span contains a figure environment, a graphics-inclusion command, a PDF-page inclusion, a table or a TikZ picture; no image file is copied into this package, the package declares no asset dependency and no image file is redistributed with it. Licence: version arXiv:2310.07778v3 of this article is distributed under the Creative Commons Attribution 4.0 International licence, as recorded in the arXiv licence metadata for this exact version and shown on the abstract page https://arxiv.org/abs/2310.07778v3. A full rights notice is delivered at licenses/arxiv-2310.07778-CC-BY-4.0.txt. Characters: every retained excerpt is pure ASCII, so the delivered engine, XeLaTeX with its default Latin Modern text font, reports no missing character.
\begin{Lem}\label{et1}
The morphism between quotient stacks 
\begin{equation}
[S[n]/G[n]]\longrightarrow \mathfrak M
\end{equation}
is \'{e}tale.
\end{Lem}

\begin{Prop}\label{rellog101}
The morphism $\mathfrak M\longrightarrow S$ is a log-smooth map. Moreover, the relative log-cotangent complex $\mathbb L^{\log}_{\mathfrak M/S}=0$. 
\end{Prop}

\begin{proof}
From Lemma \ref{et1}, it follows that the map $[S[n]/ G[n]]\longrightarrow \mathfrak M$ is \'{e}tale. Notice that the varieties $S[n]$ and $S$ are log smooth varieties with the log structure coming from the divisor given by the pre-image of the closed point of $S$ under the map $S[n]\longrightarrow S$. Since the map $S[n]\longrightarrow S$ is an \'{e}tale base-change of the map $\mathbb A^{n+1}\longrightarrow \mathbb A^1$ given by $(t_1, \dots, t_{n+1})\mapsto t_1\cdots t_{n+1}$, the map is a log-smooth morphism. Therefore, it follows that the map $[S[n]/ G[n]]\longrightarrow S$ is also a log-smooth morphism of log-smooth Artin stacks. 

The explicit description of the relative log-cotangent complex is the following. 
The pull-back (to $S[n]$) of the log-cotangent complex of the stack $[S[n]/G[n]]$ is the perfect 
complex on $S[n]$ concentrated in degrees $0$ and $1$ which is given by 
\begin{equation}\label{moment}
\left[\Omega_{S[n]}^{1}(\log ~\partial S[n])\longrightarrow \mathcal O_{S[n]}^{\oplus n}\right]
\end{equation} 
where $\Omega_{S[n]}^{1}(\log ~\partial S[n])$ is the rank $(n+1)$ locally free sheaf of logarithmic one forms on $S[n]$ with poles along the strict normal crossings divisor $\partial S[n] = S[n]\times_{S} \{o\}$.  Since the action of $G[n]$ on $S[n]$ respects $\partial S[n]$, we obtain a natural $G[n]$-equivariant structure $\Omega_{S[n]}^{1}(\log ~\partial S[n])$  such that the natural log-symplectic structure is $G[n]$-equivariant. The morphism \eqref{moment} is given by the (equivariant)  moment map for the action of $G[n]$ on the log-cotangent bundle of $S[n]$. In particular \eqref{moment} is a $G[n]$-perfect complex. We recall that the action of $G[n]$ on $S[n]$ is given by 
\begin{equation}\label{act123}
(t_1, \dots, t_{n+1})\cdot(x_1, \dots, x_{n+1})=(t_1x_1, t_2 x_2, \dots, t_{n+1}x_{n+1})
\end{equation}
where $t_1\cdots t_{n+1}=1$, and the map $S[n]\longrightarrow S$ is given by

\begin{equation}\label{act}
(t_1, \dots, t_{n+1})\mapsto \prod^{n+1}_{i=1} t_i
\end{equation}

It is easy to see that the map \eqref{moment} 
\begin{equation}
\left[\Omega_{S[n]}^{1}(\log ~\partial S[n])\longrightarrow \mathcal O_{S[n]}^{\oplus n}\right]
\end{equation} 
of vector bundles is surjective: consider the action of $(\mathbb G_m)^n$ on $\mathbb A^{n+1}$ given by \eqref{act123}. Let $p:=(x_1, \cdots, x_{n+1})\in \mathbb A^{n+1}$ be any point. Then consider the orbit map at the point $p$. 

\begin{align}
    Or_p: (\mathbb G_m)^n\longrightarrow \mathbb A^{n+1}\\
    (t_1, \cdots, t_{n+1})\mapsto (t_1\cdot x_1, \cdots, t_{n+1}\cdot x_{n+1})
\end{align}

The induced map on the tangent spaces is $d(Or)_p: T_{(1,\cdots, 1)}(\mathbb G_m)^n=\mathbb k^n\longrightarrow T_p \mathbb A^{n+1}=\mathbb k^{n+1}$ is given by the multiplication by the following diagonal matrix

\smallskip

\[
\begin{bmatrix}
    x_1       & 0 & 0 & \dots & 0 \\
    0       & x_2 & 0 & \dots & 0 \\
    \hdotsfor{5} \\
    0       & 0 & 0 & \dots & x_{n+1}
\end{bmatrix}
\]

\smallskip

This implies that the image of $[\sum^{n+1}_{i=1} \lambda_i \partial_{t_i}]\in T_{(1,\cdots, 1)}(\mathbb G_m)^n=\mathbb k^n$ (notice that $\sum^{n+1}_{i=1} \lambda_i=0$) is $\sum^{n+1}_{i=1} \lambda_i\cdot x_i \cdot \partial_{x_i}=\sum^{n+1}_{i=1} \lambda_i\cdot \partial_{\log~x_i}$. Therefore, we get a map 

\begin{align}
    d(Or)_p: T_{(1,\cdots, 1)}(\mathbb G_m)^n=\mathbb k^n\longrightarrow (T \mathbb A^{n+1}(-\log~\partial A^{n+1}))_p=\mathbb k^{n+1}\\
    (\sum^{n+1}_{i=1} \lambda_i\cdot \partial_{t_i})\mapsto \sum^{n+1}_{i=1} \lambda_i\cdot \partial_{\log~x_i}
\end{align}

The map above is clearly injective. Therefore, the dual of this map 

\begin{align}
    (d(Or)_p)^{\vee}: \Omega^1_p \mathbb A^{n+1}(\log~\partial A^{n+1})\longrightarrow \Omega^1_{(1,\cdots, 1)}(\mathbb G_m)^n
\end{align}

is surjective. This is the explicit description of the moment map \ref{moment}. It is easy to see that the composite map
\begin{equation}
\Omega_{S}^{1}(\log o)\hookrightarrow \Omega_{S[n]}^{1}(\log ~~\partial S[n])\longrightarrow \mathcal O_{S[n]}^{\oplus n}
\end{equation}
is $0$, and that the morphism $\Omega_{S[n]/S}^{1}(\log ~~\partial S[n])\longrightarrow \mathcal O_{S[n]}^{\oplus n}$ is an isomorphism.\\

Notice that the perfect complex
\begin{equation}
\big[\Omega_{S[n]/S}^{1}(\log ~~\partial S[n])\longrightarrow \mathcal O_{S[n]}^{\oplus n}\big]
\end{equation}
is precisely the relative log cotangent complex of the morphism $\big[ S[n]/G[n]\big]\longrightarrow S$ pulled back to $S[n]$. But since the morphism $\Omega_{S[n]/S}^{1}(\log ~~\partial S[n])\longrightarrow \mathcal O_{S[n]}^{\oplus n}$ is an isomorphism, the relative log-cotangent complex of the morphism $\big[ S[n]/G[n]\big]\longrightarrow S$ is equivalent to $0$.
\end{proof}

\end{document}
